\documentclass[aspectratio=169]{beamer}

\usetheme{Madrid}
\usecolortheme{default}
\setbeamertemplate{navigation symbols}{}

\usepackage{graphicx}
\usepackage{booktabs}
\usepackage{tabularx}
\usepackage{array}

\definecolor{primaryblue}{RGB}{26, 77, 140}
\definecolor{accentgold}{RGB}{199, 138, 32}
\definecolor{softgray}{RGB}{82, 82, 82}

\setbeamercolor{title}{fg=primaryblue}
\setbeamercolor{frametitle}{fg=primaryblue}
\setbeamercolor{structure}{fg=primaryblue}
\setbeamercolor{block title}{bg=primaryblue!12, fg=primaryblue}
\setbeamercolor{block body}{bg=primaryblue!4, fg=black}
\setbeamercolor{footline}{bg=primaryblue, fg=white}

\setbeamerfont{title}{series=\bfseries,size=\LARGE}
\setbeamerfont{frametitle}{series=\bfseries,size=\large}

\title[Secure Digital Evidence Platform]{Secure Digital Evidence Management\\and Forensics Training Platform}
\subtitle{Project Presentation}
\author{Neha Gaikwad}
\institute{School of Engineering\\Ajeenkya DY Patil University, Pune}
\date{April 2026}

\logo{\includegraphics[height=0.8cm]{horizontal college logo.png}}

\begin{document}

\begin{frame}[plain]
    \titlepage
\end{frame}

\section{Introduction}

\begin{frame}{Background and Motivation}
    \begin{itemize}
        \item Modern organizations rely heavily on digital documents, images, reports, and case files.
        \item Sensitive files used in forensics and legal workflows require confidentiality, controlled access, and accountability.
        \item Ordinary file-sharing systems do not preserve evidence handling history in a meaningful way.
        \item Academic projects often explain security concepts separately but do not integrate them into one practical workflow.
    \end{itemize}
\end{frame}

\begin{frame}{Problem Statement}
    \begin{block}{Core Problem}
        Traditional file handling systems may leave important evidence unencrypted, insufficiently tracked, and vulnerable to unauthorized access or misuse.
    \end{block}

    \vspace{0.3cm}

    \begin{itemize}
        \item No structured chain-of-custody awareness
        \item Weak access control for sensitive files
        \item Limited visibility into who performed evidence actions
        \item Lack of an educational platform for secure evidence handling practice
    \end{itemize}
\end{frame}

\begin{frame}{Objectives, Scope, and End Users}
    \begin{columns}[T]
        \column{0.5\textwidth}
        \textbf{Objectives}
        \begin{itemize}
            \item Secure web platform
            \item JWT-based access
            \item AES-256-CBC protection
            \item Metadata preservation
            \item Inspection and logging
        \end{itemize}

        \column{0.46\textwidth}
        \textbf{Scope and Users}
        \begin{itemize}
            \item Encryption and decryption
            \item Metadata inspection
            \item Evidence activity log
            \item WebSocket guidance
            \item Trainees
            \item Supervisors
            \item Trainers
        \end{itemize}
    \end{columns}
\end{frame}

\section{Concept and Design}

\begin{frame}{Literature and Design Insights}
    \begin{itemize}
        \item Digital evidence systems must preserve integrity, confidentiality, and accountability.
        \item AES is widely accepted for strong and efficient data protection.
        \item Chain of custody is essential in digital forensics because it documents evidence handling history.
        \item Secure web applications commonly rely on token-based authentication and protected routes.
    \end{itemize}

    \begin{block}{Design Decision}
        The project combines these ideas into a simplified training platform that is realistic enough for demonstration but manageable in an academic setting.
    \end{block}
\end{frame}

\begin{frame}{Proposed Solution}
    \begin{block}{What the Platform Does}
        The system allows authenticated users to encrypt files into evidence packages, inspect package details, decrypt evidence when needed, and review a log of evidence actions.
    \end{block}

    \begin{itemize}
        \item Protects file contents with AES-256-CBC
        \item Controls access with JWT-based authentication
        \item Preserves forensic metadata with each package
        \item Supports custody-aware inspection before decryption
        \item Records evidence actions for audit visibility
    \end{itemize}
\end{frame}

\begin{frame}{System Architecture}
    \begin{center}
    \renewcommand{\arraystretch}{1.4}
    \begin{tabular}{>{\centering\arraybackslash}m{0.24\textwidth}>{\centering\arraybackslash}m{0.08\textwidth}>{\centering\arraybackslash}m{0.28\textwidth}>{\centering\arraybackslash}m{0.08\textwidth}>{\centering\arraybackslash}m{0.22\textwidth}}
        \textbf{Frontend}\\
        HTML, CSS, JavaScript
        &
        $\rightarrow$
        &
        \textbf{Express Server}\\
        Auth, encryption, decryption, inspection, logging
        &
        $\rightarrow$
        &
        \textbf{Outputs}\\
        Encrypted package, decrypted file, metadata, log
    \end{tabular}
    \end{center}

    \vspace{0.4cm}

    \begin{itemize}
        \item Client-server architecture
        \item Browser communicates through HTTP requests
        \item WebSocket channel supports training guidance updates
    \end{itemize}
\end{frame}

\begin{frame}{Core Modules and Workflow}
    \begin{columns}[T]
        \column{0.48\textwidth}
        \textbf{Major Modules}
        \begin{itemize}
            \item Authentication module
            \item Evidence encryption module
            \item Evidence decryption module
            \item Evidence inspection module
        \end{itemize}

        \column{0.48\textwidth}
        \begin{itemize}
            \item Evidence activity log module
            \item WebSocket guidance module
            \item User interface dashboard
        \end{itemize}
    \end{columns}

    \vspace{0.2cm}

    \begin{block}{Implementation Files}
        \texttt{src/server.js}, \texttt{public/index.html}, \texttt{public/style.css}, \texttt{public/websocket\_client.js}
    \end{block}
\vspace{0.1cm}
\textbf{Workflow:} Login $\rightarrow$ Token issued $\rightarrow$ Upload file $\rightarrow$ Encrypt package $\rightarrow$ Inspect metadata $\rightarrow$ Decrypt if authorized $\rightarrow$ Log action
\end{frame}

\section{Implementation}

\begin{frame}{Technology Stack}
    \begin{table}
        \centering
        \begin{tabularx}{\textwidth}{>{\bfseries}l X}
            \toprule
            Layer & Technology \\
            \midrule
            Frontend & HTML, CSS, JavaScript \\
            Backend & Node.js, Express.js \\
            Security & JWT authentication, AES-256-CBC encryption \\
            File Handling & Multer memory storage, Node.js Buffer processing \\
            Real-Time Support & WebSocket server \\
            \bottomrule
        \end{tabularx}
    \end{table}
\end{frame}

\begin{frame}{Security Design}
    \begin{itemize}
        \item JWT tokens protect evidence-related routes after login.
        \item AES-256-CBC encrypts file contents before storage or transfer.
        \item Every encrypted package includes forensic metadata such as:
        \begin{itemize}
            \item original file name
            \item MIME type
            \item file size
            \item evidence ID
            \item encryption timestamp
            \item chain-of-custody events
        \end{itemize}
        \item Evidence actions are recorded with timestamps and short hashes.
    \end{itemize}
\end{frame}

\begin{frame}{Evidence Package Structure}
    \begin{block}{Packaging Logic}
        Before encryption, the system combines:
    \end{block}

    \begin{enumerate}
        \item metadata length
        \item metadata JSON
        \item original file content
    \end{enumerate}

    \vspace{0.3cm}

    \begin{block}{Why This Matters}
        This structure allows the platform to decrypt the package later, recover metadata, inspect details, and maintain simplified chain-of-custody awareness.
    \end{block}
\end{frame}

\begin{frame}{Major API Endpoints}
    \begin{table}
        \centering
        \begin{tabularx}{\textwidth}{>{\ttfamily}l X}
            \toprule
            Route & Purpose \\
            \midrule
            POST /api/authenticate & Validate credentials and issue JWT token \\
            POST /api/encrypt-evidence & Encrypt uploaded evidence and generate package \\
            POST /api/decrypt-evidence & Decrypt a valid evidence package \\
            POST /api/inspect-evidence & Show package metadata before decryption \\
            GET /api/evidence-log & Retrieve activity history for evidence actions \\
            \bottomrule
        \end{tabularx}
    \end{table}
\end{frame}

\begin{frame}{User Interface Highlights}
    \begin{itemize}
        \item Clean login page for secure training access
        \item Dashboard sections for:
        \begin{itemize}
            \item Guidance Feed
            \item Secure Evidence
            \item Access Evidence
            \item Evidence Inspector
            \item Evidence Activity
        \end{itemize}
        \item Minimal layout keeps the demonstration easy to understand
        \item Role-aware welcome flow improves training context
    \end{itemize}
\end{frame}

\section{Evaluation}

\begin{frame}{Testing and Validation}
    \begin{table}
        \centering
        \small
        \begin{tabularx}{\textwidth}{|X|X|}
            \hline
            \textbf{Test Case} & \textbf{Expected Result} \\ \hline
            Valid login & Token generated and dashboard displayed \\ \hline
            Invalid login & Error message displayed \\ \hline
            Encrypt valid file & Downloadable \texttt{.enc} file generated \\ \hline
            Decrypt valid package & Original file returned \\ \hline
            Inspect valid package & Metadata displayed successfully \\ \hline
            View evidence log & Activity records displayed \\ \hline
        \end{tabularx}
    \end{table}
\end{frame}

\begin{frame}{Key Results and Discussion}
    \begin{itemize}
        \item The platform successfully integrates authentication, encryption, metadata handling, inspection, and logging.
        \item The Evidence Inspector adds a realistic forensic verification step before decryption.
        \item The dashboard-based UI improves clarity during demonstration and project review.
        \item The project shows that cybersecurity concepts can be taught effectively through a working workflow, not just theory.
    \end{itemize}
\end{frame}

\begin{frame}{Sample Demonstration Flow}
    \begin{enumerate}
        \item Log in using a valid training account.
        \item Upload a sample PDF, image, or text file.
        \item Download the generated encrypted \texttt{.enc} package.
        \item Inspect the package to view evidence ID, file details, and custody events.
        \item Decrypt the same package and verify original content recovery.
        \item Open the evidence activity log to review recorded actions.
    \end{enumerate}
\end{frame}

\begin{frame}{Strengths of the Project}
    \begin{itemize}
        \item Clear end-to-end secure evidence workflow
        \item Easy to demonstrate in academic presentations
        \item Combines web development with cybersecurity concepts
        \item Includes custody-style logging for evidence awareness
        \item Modular structure supports future enhancement
    \end{itemize}
\end{frame}

\begin{frame}{Current Limitations}
    \begin{itemize}
        \item Activity log is not yet backed by persistent database storage.
        \item Role separation is present in training accounts but not fully enforced as advanced RBAC.
        \item Encryption key should be fully managed through environment configuration.
        \item The platform is a prototype and not a legal-grade production evidence system.
        \item Mobile responsiveness and report export can be improved further.
    \end{itemize}
\end{frame}

\section{Conclusion}

\begin{frame}{Future Enhancements}
    \begin{itemize}
        \item Add persistent database storage for evidence records
        \item Implement stronger role-based authorization
        \item Generate exportable reports for evidence activity
        \item Improve notifications and dashboard usability
        \item Prepare the system for cloud hosting
        \item Optimize the interface for mobile devices
        \item Move all secrets and keys into environment variables
        \item Add search and filtering for evidence logs
    \end{itemize}
\end{frame}

\begin{frame}{Conclusion}
    \begin{block}{Summary}
        The project delivers a practical and academic-friendly platform for secure digital evidence management by combining authentication, encryption, inspection, decryption, and activity logging in one coherent workflow.
    \end{block}

    \vspace{0.3cm}

    \begin{itemize}
        \item Demonstrates real security concepts through implementation
        \item Supports forensic awareness and evidence accountability
        \item Provides a solid foundation for future expansion
    \end{itemize}
\end{frame}

\begin{frame}[plain]
    \begin{center}
        \vspace{1.2cm}
        {\Huge \textbf{Thank You}}\\[0.5cm]
        {\Large Questions and Discussion}

        \vspace{1cm}
        \includegraphics[height=2.3cm]{horizontal college logo.png}

        \vfill
        {\large Neha Gaikwad}\\
        {\large B.Tech Information Technology}
    \end{center}
\end{frame}

\end{document}
